%% kapa: exam=1 round=26 year=2015 period=1 form=A subject=economics subject_ko=경제원론 seq=28 number=68
%% source: Q-Net 감정평가사 기출문제 (cst00309) artlSeq=5138001 file=extracted/1008/5138001/제26회_감정평가사 1차_기출/제26회_감정평가사 1차_기출/1교시_감평_A형.pdf
%% conversion: kordoc 4.24.51 -> tools/gen1.py (auto); review=proofread-sonnet; flags=none
\begin{question}{68}
다음 중 옳은 것을 모두 고른 것은? (단, 피셔효과가 성립한다.)
\begin{framedtext}
ㄱ. 실질이자율은 명목이자율에서 인플레이션율을 뺀 것이다. \par
ㄴ. 예상보다 높은 인플레이션율은 채무자에게 유리하고 채권자에게는 불리하다. \par
ㄷ. 예상되는 미래인플레이션율의 상승은 예상되는 실질이자율을 상승시킨다. \par
\end{framedtext}
\begin{choices}
  \item ㄱ
  \item ㄴ
  \item ㄱ, ㄴ
  \item ㄱ, ㄷ
  \item ㄴ, ㄷ
\end{choices}
\end{question}
\answer{3}
